% Generated by Listas/make_statements.py from pratica-lecture-07.md, item 11. Do not edit.
\textbf{Source.} Tadelis (2013), \emph{Game Theory: An Introduction}, Exercise 12.3 ``Armed Conflict'', pp.~266--267 {[}PDF pp.~285--286{]}. Reproduced verbatim. \textbf{Answer in the source:} no (none in the book; Princeton UP's website has about half, not checked).



\begin{quote}

\textbf{Armed Conflict:} Consider the following strategic situation: Two rival armies plan to seize a disputed territory. Each army's general can choose either to attack (\(A\)) or to not attack (\(N\)). In addition, each army is either strong (\(S\)) or weak (\(W\)) with equal probability, and the realizations for each army are \emph{independent}. Furthermore the type of each army is known only to that army's general. An army can capture the territory if either (i) it attacks and its rival does not or (ii) it and its rival attack, but it is strong and the rival is weak. If both attack and are of equal strength then neither captures the territory. As for payoffs, the territory is worth \(m\) if captured and each army has a cost of fighting equal to \(s\) if it is strong and \(w\) if it is weak, where \(s < w\). If an army attacks but its rival does not, no costs are borne by either side. Identify all the pure-strategy Bayesian Nash equilibria of this game for the following two cases, and briefly describe the intuition for your results:



\begin{enumerate}

\def\labelenumi{\alph{enumi}.}

\item

  \(m = 3\), \(w = 2\), \(s = 1\).

\item

  \(m = 3\), \(w = 4\), \(s = 2\).

\end{enumerate}

\end{quote}

